\documentclass{article}

% CoRL 2026 workshop style. Omit [final] for the double-blind submission
% (this anonymizes the author block); add [final] for a camera-ready version.
% Download corl_2026.sty from the CoRL 2026 author instructions and place it
% alongside this file (or use the official CoRL 2026 Overleaf template).
\usepackage{corl_2026}

\usepackage[T1]{fontenc}
\usepackage[utf8]{inputenc}
\usepackage{microtype}
\usepackage{amsmath,amssymb}
\usepackage{booktabs}
\usepackage{tabularx}
\usepackage{array}
\usepackage{xcolor}
\usepackage{graphicx}
\usepackage{tikz}
\usetikzlibrary{positioning}
\usepackage{hyperref}

\hypersetup{colorlinks=true, linkcolor=blue, citecolor=blue, urlcolor=blue,
  pdftitle={Diffusion-Trained Backbones for Flow-Matching VLA Policies}}

\title{Diffusion-Trained Backbones for Flow-Matching\\
Vision-Language-Action Policies}

% Double-blind submission: authors withheld. Do not de-anonymize until acceptance.
\author{
  Anonymous Author(s)\\
  Affiliation withheld for review\\
  \texttt{anonymous@example.com}
}

\begin{document}
\maketitle

\begin{abstract}
Vision-Language-Action (VLA) policies pair a pretrained multimodal backbone with
a continuous action generator such as a flow-matching action expert. Most
backbones inherit an autoregressive (AR) next-token pretraining prior, yet robot
action chunks are block-structured rather than left-to-right. We ask whether a
\emph{diffusion-trained} bidirectional backbone produces better action-token
hidden states than a scale-matched AR backbone when both drive the same
flow-matching action expert. We isolate the pretraining prior with a matched pair
--- Dream-7B (discrete diffusion, initialized from Qwen2.5-7B) versus Qwen2.5-7B
(AR) --- holding architecture, initialization, action expert, dataset, and control
stack fixed, and evaluate entirely in simulation on the LIBERO benchmark with
both policy-level success and representation-level probing. On the matched $C-D$
contrast (single seed, 500 rollouts/suite), the diffusion-trained backbone improves
mean LIBERO success by $+5.9$ points (82.2 vs.\ 76.3). The gain is concentrated
exactly where the autoregressive backbone is weakest: it is large and statistically
significant (disjoint 95\% Wilson intervals, $p<0.001$) on the long-horizon
($+12.0$) and object ($+11.2$) suites, and a tie on the near-ceiling spatial and
goal suites. Linear probing of the two backbones' action-token features, however,
finds them equally decodable for low-level action, state, and gripper targets
($|C-D|\le0.01$): the behavioral advantage is not reducible to the linear
decodability of immediate control signals, which rules out the simplest
explanation and leaves the underlying mechanism to future probing. We report the study as a controlled, single-seed result and
discuss its limitations, including seed variance and the objective-vs-corpus
confound inherent to off-the-shelf backbones.
\end{abstract}

\keywords{Vision-Language-Action Policies, Flow Matching, Diffusion Models,
Imitation Learning, Representation Probing}

\section{Introduction}
State-of-the-art vision-language-action (VLA) policies pair a pretrained
vision-language model (VLM) backbone with a flow-matching action expert that turns
the backbone's hidden states into continuous action chunks, as in $\pi_0$ and
$\pi_{0.5}$ \cite{pi0,pi05,openpi}. The action expert generates a whole block of
future actions at once, refining them iteratively and attending across the block
bidirectionally --- a generative process much closer to denoising than to
left-to-right prediction. The backbone that feeds it, however, is almost always an
\emph{autoregressive} VLM, pretrained to predict the next token causally. This
raises a question that the field has largely left implicit: does the backbone's
\emph{pretraining objective} --- the prior it brings to the action-token hidden
states --- matter for the policy, and would a backbone pretrained by denoising
suit a flow-matching expert better than a matched autoregressive one?

The question has been hard to ask cleanly for two reasons. First, until recently
there was no diffusion-pretrained language model at the scale of the autoregressive
backbones used in VLAs; discrete-diffusion LMs such as Dream-7B
\cite{dream}, initialized from Qwen2.5-7B \cite{qwen25}, now provide an
architecturally identical counterpart whose defining difference is a
discrete-diffusion pretraining objective. Second, the inference-time attention mask is easy to confuse with the
pretraining prior: we note that the $\pi_{0.5}$ backbone is in fact a prefix-LM
whose action block is already attended bidirectionally, so the causal-vs-bidirectional
distinction at inference is not the variable of interest --- the pretraining prior
is. These two observations make a controlled comparison possible: hold the action
expert, vision encoder, data, and adaptation recipe fixed, and swap only the
backbone between a diffusion-trained model (Dream-7B) and its autoregressive
initializer (Qwen2.5-7B).

We run exactly this comparison on LIBERO. The diffusion-trained backbone yields
the better policy --- $+5.9$ points of mean success and $+12.0$ on the
long-horizon suite, precisely where the autoregressive backbone is weakest --- at
matched scale, initialization, and everything downstream of the backbone, so the
gap is attributable to the pretraining prior. Probing the two backbones'
action-token features, however, finds them equally linearly decodable for
low-level control targets: the behavioral advantage is real but does not reduce to
the linear readability of immediate actions --- ruling out the simplest
explanation, even if our low-level probes cannot themselves localize where the
advantage lies.

Our contributions are: (i) a scale/architecture-matched isolation of the
diffusion pretraining prior for VLA backbones; (ii) a precise characterization of
the \(\pi_{0.5}\) backbone as a prefix-LM (bidirectional action block), separating
attention-mask topology from the pretraining prior; (iii) representation-level
probing of action-token hidden states; and (iv) a controlled LIBERO evaluation
with rollout-level statistics (Wilson intervals, two-proportion tests) and an
explicit account of the study's confounds and single-seed scope.

\section{Related Work}

\paragraph{Diffusion and flow policies.} Generating actions by denoising or flow
matching is now standard in imitation learning: Diffusion Policy \cite{diffusion_policy}
and 3D Diffusion Policy \cite{dp3} show that iterative, block-structured action
generation captures multimodal behavior better than regressing a single action.
These methods build the denoiser on task-specific or from-scratch encoders; we
inherit the flow-matching action head but ask what its \emph{backbone} should be.

\paragraph{VLM backbones with flow-matching experts.} $\pi_0$ and $\pi_{0.5}$
\cite{pi0,pi05,openpi} scale this recipe by conditioning a flow-matching expert on
a large pretrained VLM, achieving strong generalist manipulation. In these systems
the backbone is autoregressively pretrained and (in $\pi_{0.5}$) used as a
prefix-LM, so the action block is already bidirectional at inference; the
pretraining objective of the backbone is not itself treated as a design choice.
Our study takes that objective as the variable.

\paragraph{Diffusion-trained and multimodal-diffusion LMs.} Discrete-diffusion
language models --- Dream-7B \cite{dream}, and encoder-decoder diffusion LMs such
as DiffusionGemma \cite{diffusiongemma_docs} --- and multimodal diffusion models
like LaViDa \cite{lavida} establish that denoising-pretrained backbones exist at
VLM scale. Crucially, Dream is initialized from an autoregressive model
(Qwen2.5-7B), giving us a matched pair whose principal difference is the
pretraining objective (with the corpus/compute caveat we bound in
\S\ref{sec:limits}).

\paragraph{Unified-backbone diffusion VLAs (closest work).} A recent line folds
action decoding \emph{into} a diffusion backbone --- Dream-VLA \cite{dreamvla},
LLaDA-VLA \cite{lladavla}, and Discrete Diffusion VLA \cite{discretediffvla} ---
reporting strong LIBERO results. These works differ from ours in three ways that
matter for what can be concluded: they (a) do not hold a matched autoregressive
backbone fixed as a control, so the effect of the pretraining prior is
confounded with architecture and training changes; (b) let the backbone generate
actions directly rather than through a separate, fixed flow-matching expert, so
representation quality and decoding are entangled; and (c) do not probe the
learned representations. We instead keep one external flow-matching expert
constant, swap only the backbone against its own AR initializer, and add
representation-level probing --- trading peak performance for a controlled,
interpretable contrast.

\paragraph{Representation quality and action-token structure.} Crossway Diffusion
\cite{crossway} shows that auxiliary representation learning improves diffusion
policies, and work on ordered/structured action tokens \cite{oat} highlights that
how actions are laid out as tokens affects learning --- both motivating our focus
on the \emph{action-token hidden states} the expert consumes, and our attempt to
measure their quality directly via probing rather than only through rollouts.

\section{Method}
\subsection{Backbone-swap with a shared flow-matching expert}

\begin{figure}[t]
\centering
\begin{tikzpicture}[
  font=\footnotesize,
  io/.style ={draw, rounded corners, fill=gray!8, align=center,
              text width=26mm, minimum height=6mm, inner sep=2pt},
  bb/.style ={draw=orange!75!black, thick, rounded corners,
              fill=orange!15, align=center, text width=26mm, minimum height=27mm},
  exp/.style={draw=blue!65!black, thick, rounded corners,
              fill=blue!12, align=center, text width=24mm, minimum height=17mm},
  out/.style={draw, rounded corners, fill=gray!8, align=center,
              text width=15mm, minimum height=10mm},
]
% inputs (stacked)
\node[io] (img)  {Image $\to$ SigLIP + connector};
\node[io, below=2.5mm of img]   (lang)  {Language $\to$ embed};
\node[io, below=2.5mm of lang]  (state) {Proprioceptive state};
\node[io, below=2.5mm of state] (act)   {Flow-time $\tau$ + noisy actions};
% backbone (the swap)
\node[bb, right=11mm of lang.east, yshift=-7mm, anchor=west] (bb)
  {\textbf{Backbone} \\ {\scriptsize\itshape only part that changes} \\[3pt]
   \textbf{C:} Dream-7B \\ {\scriptsize discrete diffusion} \\[2pt]
   \textbf{D:} Qwen2.5-7B \\ {\scriptsize autoregressive}};
% shared standalone expert
\node[exp, right=12mm of bb, anchor=west] (exp)
  {\textbf{Standalone flow expert} \\[2pt] {\scriptsize DiT --- identical for C, D}};
% output
\node[out, right=10mm of exp, anchor=west] (out) {Action chunk $a_{1:H}$};
% arrows
\draw[->] (img.east)   -- (bb.west);
\draw[->] (lang.east)  -- (bb.west);
\draw[->] (state.east) -- (bb.west);
\draw[->] (act.east)   -- (bb.west);
\draw[->] (bb.east)  -- node[above, font=\scriptsize] {$H_{\text{action}}$} (exp.west);
\draw[->] (exp.east) -- node[above, font=\scriptsize] {$v_\theta$} (out.west);
\end{tikzpicture}
\caption{\textbf{Controlled backbone swap.} Vision (SigLIP + connector), the
standalone flow-matching action expert, the action space, dataset, adaptation
recipe, and the block-diagonal attention mask are held \emph{identical}; only the
language backbone is swapped between the diffusion-trained Dream-7B (C) and its
autoregressive initializer Qwen2.5-7B (D). The expert reads the action-token
hidden states $H_{\mathrm{action}}$ and outputs the flow vector field $v_\theta$
over the action chunk. Because everything but the backbone is matched, the $C-D$
difference isolates the pretraining prior.}
\label{fig:method}
\end{figure}

Image/language/proprio/flow-time/noisy-action tokens are processed by the
backbone, which contextualizes the action tokens (Figure~\ref{fig:method}); the hidden states at the
action-token positions, $H_{\text{action}}$, are read by a flow-matching action
expert that outputs the vector field over the action chunk. The flow objective
is $\mathcal{L}=\lVert v_\theta - (A_{\text{noise}}-A_{\text{clean}})\rVert^2$
with $x_t = t\,A_{\text{noise}} + (1-t)\,A_{\text{clean}}$.

\paragraph{Standalone vs.\ fused expert.} A subtlety determines the design.
$\pi_{0.5}$'s action expert is not a detachable head: it is a second transformer
stream \emph{fused} into the backbone's self-attention (per layer, backbone and
expert queries/keys/values are concatenated and attended jointly, sharing RoPE
and KV cache), which forces the expert to match the backbone's layer count, head
geometry, and normalization. This fusion cannot be preserved when the backbone
changes from Gemma to Qwen2.5/Dream (different depth, heads, RoPE, biases). We
therefore use a \emph{standalone} flow expert --- a small DiT-style transformer
over the action tokens, flow-time injected via adaLN-Zero conditioning --- that
reads $H_{\text{action}}$ and is decoupled from the backbone's attention. It is
shared \emph{identically} across variants A$'$/C/D; the only variant-specific
parameter is a linear input adapter mapping the backbone width to the common
expert width ($\sim$2--4M of $\sim$117M parameters), so the C$-$D contrast
remains clean. Model A (stock $\pi_{0.5}$, fused expert) is retained as a
reference that reproduces published numbers, and A$'$ (Gemma backbone $+$ the
standalone expert) is the architecture-matched baseline.

\subsection{Model variants}
The action expert, action space, flow objective, dataset, and control stack are
identical across all rows; only the backbone/attention topology changes.

\begin{table}[t]
\centering
\caption{Model variants. Identifying contrast: $C-D$ (diffusion prior at matched
scale/architecture).}
\begin{tabularx}{\linewidth}{>{\bfseries}p{0.06\linewidth} X}
\toprule
 & Backbone \\ \midrule
A & Stock $\pi_{0.5}$ (Gemma, \emph{fused} expert) --- reference baseline \\
A$'$ & Gemma backbone $+$ standalone expert --- architecture-matched baseline \\
D & Qwen2.5-7B (AR) $+$ standalone expert --- matched control \\
C & Dream-7B (diffusion, init.\ from Qwen2.5-7B) $+$ standalone expert --- proposed \\
\bottomrule
\end{tabularx}
\end{table}

\subsection{Reading hidden states from a diffusion backbone}
\label{sec:kread}
A diffusion-trained backbone could in principle be run for several internal
refinement steps ($k>1$) before its hidden states are read, unrolling the
denoising process the AR backbone lacks. To keep the $C-D$ contrast about the
pretraining \emph{prior} and not about extra inference compute, we read both
backbones with a single forward pass ($k{=}1$): C and D then do exactly the same
amount of work per action, and any gap is due to the learned representation, not
to additional refinement. This backbone read is also distinct from --- and
happens once per --- each flow-matching step of the standalone expert. We report
$k{=}1$ throughout; multi-step refinement reads ($k\in\{2,4\}$), which trade
compute for a potentially richer read, are a natural ablation left to future work.

\section{Experimental Setup}
Evaluation is entirely in simulation on LIBERO \cite{libero} (single-arm Franka,
7-DoF), across four suites (Spatial, Object, Goal, Long/libero\_10), 10 tasks per
suite $\times$ 50 trials $=$ 500 rollouts/suite. Backbones are $\sim$7B (Dream-7B
\cite{dream} / Qwen2.5-7B \cite{qwen25}); Model A is the stock $\pi_{0.5}$ (Gemma)
backbone. \textbf{Vision is held identical across the standalone-expert variants}
(A$'$/C/D): the SigLIP tower from PaliGemma plus a fresh linear connector to the
backbone width, so only the language backbone changes between C and D. Training
forks the openpi \texttt{pi05\_libero} recipe (30k steps, batch 32, LR $5\times
10^{-5}$, action horizon $H{=}10$, action dim padded to 32). \textbf{Adaptation
differs by variant and is stated with each result:} A and A$'$ are fully
fine-tuned from \texttt{pi05\_base}; C and D use an identical strengthened-LoRA
recipe (rank 64, $\alpha{=}128$ on the frozen 7B backbone, with the SigLIP tower,
connector, input adapters, and standalone expert trainable) --- chosen because
full fine-tuning of a 7B backbone with AdamW does not fit one 80GB GPU under DDP,
and because freezing the base best preserves the pretraining prior the $C-D$
contrast is designed to isolate. The diffusion backbone is read with a single
forward pass ($k{=}1$), matching D exactly; deeper $k$-step refinement reads are
left to future work. Each training run used 2$\times$A100 ($\sim$5\,s/step). We
report per-suite success rate with 95\% Wilson confidence intervals (500
rollouts/suite) and a two-proportion test on the $C-D$ contrast, with LIBERO-Long
as the primary discriminator, plus representation probes (\S\ref{sec:probing}). A
and A$'$ are averaged over 3 seeds; C and D are single-seed ($n{=}1$) --- the
identifying $C-D$ contrast is evaluated on the matched seed, with additional seeds
for error bars left to future work.

\paragraph{Reproducibility.} The implementation extends OpenPI \cite{openpi}
(Apache-2.0); the backbones (Dream-7B, Qwen2.5-7B) and the LIBERO benchmark are
publicly available, and no private data is used. The standalone flow expert,
per-variant training configs (LoRA rank 64, batch 32, 30k steps, LR $5\times
10^{-5}$, $k{=}1$), and the probing pipeline are implemented as an
\texttt{openpi.diffusion\_backbone} package; code and configurations will be
released.

\section{Results}

\subsection{Stage 1 --- Model A baseline reproduction}
Our Model A reproduces the published $\pi_{0.5}$ LIBERO results to within
$\sim$1\% on every suite (average 96.6 vs.\ 96.85), validating the training and
evaluation harness before any backbone swap. This holds despite training at
batch size 32 (vs.\ the published 256), supporting our batch choice. On the
long-horizon suite Model A slightly exceeds the published number (93.1 vs.\ 92.4).

\begin{table}[t]
\centering
\caption{Model A (stock $\pi_{0.5}$) LIBERO success rate (\%), mean $\pm$ sd over
3 seeds, vs published $\pi_{0.5}$@30k. Higher is better.}
\begin{tabular}{lccccc}
\toprule
Model & Spatial & Object & Goal & Long & Avg \\
\midrule
Published $\pi_{0.5}$@30k & 98.8 & 98.2 & 98.0 & 92.4 & 96.85 \\
A (ours, 3 seeds) & 97.7\,$\pm$\,0.6 & 98.3\,$\pm$\,0.6 & 97.1\,$\pm$\,0.5 & 93.1\,$\pm$\,1.4 & 96.6\,$\pm$\,0.2 \\
\bottomrule
\end{tabular}
\end{table}

\subsection{Main comparison: diffusion vs matched AR backbone}
\label{sec:main}
\begin{table}[t]
\centering
\caption{LIBERO success rate (\%) by variant. A/A$'$ are mean over 3 seeds;
D/C are single-seed ($n{=}1$, seed 0). The identifying contrast is $C-D$,
evaluated on the matched single seed. D/C use the strengthened LoRA recipe
(unfrozen SigLIP, rank 64); A$'$ is full fine-tuned. Superscripts on $C-D$ mark a
two-proportion test over the 500 rollouts per suite ($^{\ast\ast\ast}p<0.001$):
Object and Long are significant at the rollout level (95\% Wilson intervals
disjoint), Spatial and Goal are not. This quantifies rollout precision only, not
training-seed variance, which $n{=}1$ leaves unmeasured (see \S\ref{sec:limits}).}
\begin{tabular}{lccccc}
\toprule
Model & Spatial & Object & Goal & Long & Avg \\
\midrule
A (stock $\pi_{0.5}$, fused, 3 seeds) & 97.7 & 98.3 & 97.1 & 93.1 & 96.6 \\
A$'$ (Gemma + standalone, 3 seeds)    & 96.7 & 92.6 & 90.5 & 79.6 & 89.9 \\
D (Qwen2.5-7B, AR, n{=}1)             & 94.2 & 82.0 & 71.4 & 57.4 & 76.3 \\
C (Dream-7B, diff., n{=}1)   & 93.0 & 93.2 & 73.2 & 69.4 & 82.2 \\
\midrule
$C-D$ (diffusion prior)      & $-1.2$ & $+11.2^{\ast\ast\ast}$ & $+1.8$ & $+12.0^{\ast\ast\ast}$ & $+5.9$ \\
$A'\!-\!A$ (decoupling cost) & $-1.0$ & $-5.7$ & $-6.6$ & $-13.5$ & $-6.7$ \\
\bottomrule
\end{tabular}
\end{table}

\paragraph{The diffusion prior helps, and most where AR is weakest.} On the
identifying $C-D$ contrast (matched seed, identical standalone expert, LoRA
recipe, and vision), the diffusion-trained backbone improves mean success by
$+5.9$ points. The gain is not uniform, and the pattern is the informative result.
On the two suites where the AR backbone is strong, the backbones are
indistinguishable: spatial ($C$ 93.0 [90.4, 94.9] vs.\ $D$ 94.2 [91.8, 95.9]) and
goal ($C$ 73.2 [69.2, 76.9] vs.\ $D$ 71.4 [67.3, 75.2]) both have overlapping 95\%
Wilson intervals ($p=0.44$, $p=0.52$). On the two suites where the AR backbone is
weak, the diffusion prior wins decisively and the intervals are disjoint: object
($C$ 93.2 [90.6, 95.1] vs.\ $D$ 82.0 [78.4, 85.1], $+11.2$, $p<10^{-7}$) and
long-horizon ($C$ 69.4 [65.2, 73.3] vs.\ $D$ 57.4 [53.0, 61.7], $+12.0$,
$p<10^{-4}$). This matches the hypothesis: D (and A$'$) degrade sharply on long,
multi-stage tasks (A$'-$A is $-13.5$ on long), and it is exactly there that the
bidirectional, diffusion-trained representation recovers the most. Because the
comparison is at matched scale (7B), matched initialization (Dream is initialized
\emph{from} Qwen2.5-7B), and matched vision/expert/recipe/seed, the gap is
attributable to the pretraining that distinguishes the two backbones rather than to
capacity or architecture (we bound this attribution in \S\ref{sec:limits}). The
significance above is at the \emph{rollout} level (500 trials/suite) and states
only that these two trained models differ; it does not speak to training-seed
variance, which we could not measure at $n{=}1$ (\S\ref{sec:limits}).

\subsection{Representation-level probing}
\label{sec:probing}

The rollout comparison (\S\ref{sec:main}) tells us \emph{whether} the diffusion
prior helps end-to-end; probing asks \emph{why}, by measuring how much
action-relevant information is already linearly decodable from the backbone's
action-token hidden states before the flow expert acts. Because C and D share
architecture, width ($3584$), vision, LoRA recipe, and expert, and differ chiefly
in pretraining (\S\ref{sec:limits}), a probe gap between them at matched capacity
is attributable to the pretraining prior and mirrors the $C-D$ contrast at the
representation level.

\paragraph{Protocol.} For each trained backbone we run a single frozen forward
pass ($k{=}1$, matching D) over 4000 LIBERO observations at flow time $\tau{=}1$
(pure-noise action input, so the features reflect only the conditioning), and read
the last-layer hidden states at the $H{=}10$ action-token positions, mean-pooled
to a $3584$-d feature. C and D are extracted in the same dataset order, so they are
probed on \emph{identical} frames (the pipeline asserts this); because both
backbones are $3584$-d, one probe architecture is capacity-matched across them with
no dimensionality confound. We fit a linear probe (ridge for regression, logistic
for classification) --- our headline linear-decodability measure --- and a small
2-layer MLP as a nonlinear ceiling, each over 5 train/test splits.

\paragraph{Targets.} We decode the action-relevant quantities available directly
from the dataset: the immediate action and the full action chunk (real 7-DoF,
$R^2$), the proprioceptive state ($R^2$), and the gripper command (median-split,
accuracy/AUROC). More abstract targets --- subtask phase and target-object pose,
which would better probe the long-horizon regime --- require simulator-state
annotation and are left to future work, as is an isolation probe of the
off-the-shelf backbones before fine-tuning. The pipeline is
\texttt{scripts/extract\_probe\_features.py} and
\texttt{openpi.diffusion\_backbone.probing}.

\begin{table}[t]
\centering
\caption{Linear-probe predictivity of frozen action-token features (mean-pooled,
$\tau{=}1$), C vs.\ D, seed 0, 4000 matched frames, 5 splits. $R^2$ for
regression, accuracy for gripper. Both backbones are probed on \emph{identical}
frames. The MLP (nonlinear) probe did not converge (max-iter reached, high
variance) and is omitted as inconclusive.}
\label{tab:probing}
\begin{tabular}{lccc}
\toprule
Target & C (Dream) & D (Qwen) & $C-D$ \\
\midrule
action (next step), $R^2$      & 0.874 & 0.869 & $+0.006$ \\
action (full chunk), $R^2$     & 0.889 & 0.887 & $+0.002$ \\
proprioceptive state, $R^2$    & 0.964 & 0.954 & $+0.009$ \\
gripper (median-split), acc    & 0.969 & 0.969 & $+0.001$ \\
\bottomrule
\end{tabular}
\end{table}

\paragraph{The gain is not a low-level decodability effect.} At the representation
level, C and D are \emph{tied}: their action-token features are equally --- and
very highly --- linearly decodable for every low-level target ($|C-D|\le0.01$
throughout), whereas the rollout gap is $+5.9$ on average and $+12.0$ on
long-horizon (\S\ref{sec:main}). This lets us cleanly \emph{rule out} the simplest
explanation of the behavioral gain: it is not that the diffusion prior makes the
immediate action, state, or gripper command more linearly readable --- that
information is present to the same (high) degree in both backbones.

We are careful not to overread the null into a positive claim about \emph{where}
the advantage lives; the probe does not localize it, for three reasons. (i)
\textbf{Ceiling:} both backbones sit at $0.87$--$0.96$, leaving little headroom for
a linear probe to separate them, so this is ``no difference detectable on these
targets,'' not ``identical representations.'' (ii) \textbf{Target level:} the
targets are low-level and read from single, in-distribution frames, whereas the
rollout gap is a long-horizon, closed-loop, compounding-error effect that static
single-frame probing is not designed to capture. (iii) \textbf{Nonlinear probe:}
the MLP did not converge and is inconclusive. Higher-level or temporal structure
and closed-loop robustness remain plausible mechanisms, but adjudicating them is
exactly the job of the deferred abstract-target and closed-loop probes; here we
claim only the negative result the data supports.

\section{Discussion}

The matched $C-D$ contrast is positive: the diffusion-trained backbone yields a
better policy, and the advantage is largest on long-horizon control
($+5.9$ mean, $+12.0$ on long-horizon, both significant at the rollout level). Because $C$ and $D$ are identical up to
the pretraining objective (Dream is initialized from Qwen2.5-7B; same expert,
vision, LoRA recipe, seed, and data), we read this as evidence that a
bidirectional/denoising pretraining prior produces hidden states that a fixed
flow-matching expert can exploit better than those of a matched autoregressive
backbone --- most in the long-horizon, multi-stage regime where the AR backbone
degrades.

The probing analysis (\S\ref{sec:probing}) sharpens the claim by ruling out the
simplest account. If the diffusion prior simply made low-level control more
linearly readable, C would out-probe D; instead the two are tied at ceiling on
every low-level target (Table~\ref{tab:probing}) while C clearly out-rolls D. We
therefore claim only what the data supports: the advantage is \emph{not} a matter
of the linear decodability of immediate action/state signals --- both backbones
carry those equally --- and at this low-level, in-distribution resolution the probe
cannot say more. Whether the effect is higher-level or temporal structure, or
closed-loop robustness to compounding error (consistent with the gap being largest
on long-horizon tasks), is left open for the deferred abstract-target and
closed-loop / distribution-shift probes.

Two points of scope. In absolute terms C sits below A$'$ and A, but this reflects
the decoupled-expert and LoRA-vs-full-fine-tuning costs, which are shared by C and
D and cancel in the identifying $C-D$ contrast. And while our design controls
attention \emph{topology} by giving C and D the same mask, a natural follow-up is
a backbone with an \emph{expanded} attention topology over the action block (a
``Model~B'' built on D's weights), which would separate any benefit of richer
action-block attention from the pretraining prior; we leave this, and the deeper
probes above, to future work.

\subsection{Limitations}
\label{sec:limits}

\paragraph{Training-seed variance.} $C$ and $D$ are single-seed ($n{=}1$). The
Wilson intervals and tests in \S\ref{sec:main} establish that the two \emph{trained
models} differ at the rollout level, but they do not capture variance across
retraining. This matters most for long-horizon: across seeds, A$'$ long-horizon
success ranged 74.6--82.4 (sd $\approx$ 4.3), so single-seed suite gaps should be
read as provisional. That the two significant gaps ($+11.2$, $+12.0$) are several
times this seed spread is encouraging, but confirming them with additional seeds is
the most important next step.

\paragraph{Objective vs.\ corpus confound.} Dream-7B is initialized from
Qwen2.5-7B and then trained with a discrete-diffusion objective, but that training
also uses its own data and compute. With off-the-shelf backbones we cannot
separate the pretraining \emph{objective} from the \emph{corpus/compute} that
accompanied it; ``pretraining prior'' should therefore be read as ``the
pretraining that produced Dream rather than Qwen,'' not the diffusion objective in
isolation. The shared initialization narrows, but does not eliminate, this
confound; fully isolating the objective would require pretraining a matched pair
under one corpus, which is beyond our scope.

\paragraph{What is \emph{not} confounded.} Attention topology is controlled: $C$
and $D$ use the \emph{identical} block-diagonal prefix-LM attention mask at both
training and inference (Dream respects the passed 4D mask), so the contrast is not
explained by bidirectional-vs-causal attention \emph{usage} --- only by the
pretraining behind the weights. Scale, initialization, vision encoder, action
expert, adaptation recipe, data, and seed are likewise matched.

\paragraph{Scope.} Evaluation is simulation-only on a single benchmark (LIBERO),
so real-robot external validity is untested. Absolute performance reflects the
standalone (decoupled) expert, which costs $\sim$6--7 points vs.\ the fused expert
($A'-A$); this cost is shared by $C$ and $D$ and cancels in $C-D$, but it means the
study characterizes a controlled regime below state-of-the-art rather than a
production policy. Finally, the probes target low-level quantities that saturate at
ceiling for both backbones, so the representation-level comparison
(\S\ref{sec:probing}) is suggestive rather than conclusive.

\section{Conclusion}

We asked whether the pretraining objective of a VLA backbone --- diffusion versus
autoregressive --- matters for downstream manipulation when everything else is
held fixed, and built a controlled test to answer it: a single standalone
flow-matching action expert driven by interchangeable 7B backbones that differ
only in how they were pretrained (Dream, initialized from Qwen2.5-7B, versus
Qwen2.5-7B itself). The diffusion-trained backbone is the better policy, improving
mean LIBERO success by $+5.9$ points and by $+12.0$ on the long-horizon suite ---
the regime where the autoregressive backbone, and the decoupled expert generally,
degrade most. Because the two systems are matched in scale, initialization,
vision, action expert, adaptation recipe, data, and seed, we attribute the gap to
the pretraining prior rather than to capacity or architecture. At the same time,
linear probing shows the two backbones' action-token representations to be equally
decodable for low-level action, state, and gripper targets: the behavioral
advantage does not reduce to the linear readability of immediate control signals,
which rules out the simplest explanation while leaving the mechanism open.
The study is single-seed, simulation-only, and probes a deliberately low-level
target set, so we read it as a controlled existence proof rather than a final
measurement: a diffusion-trained backbone can yield a better flow-matching policy
than a matched autoregressive one, most where it matters, and not for the reason
one would first guess. Establishing the mechanism --- through abstract-target and
closed-loop probing --- and the generality --- through additional seeds, the
attention-topology control ($B$), and real-robot evaluation --- is the natural
next step.

\begin{thebibliography}{99}
\bibitem{diffusion_policy} Chi et al. \textit{Diffusion Policy}. RSS 2023. \url{https://arxiv.org/abs/2303.04137}
\bibitem{pi0} Physical Intelligence. \textit{$\pi_0$}. 2024. \url{https://www.physicalintelligence.company/download/pi0.pdf}
\bibitem{pi05} Black et al. \textit{$\pi_{0.5}$}. CoRL 2025. \url{https://proceedings.mlr.press/v305/black25a.html}
\bibitem{openpi} Physical Intelligence. \textit{OpenPI}. \url{https://github.com/Physical-Intelligence/openpi}
\bibitem{dream} Ye, Xie et al. \textit{Dream 7B}. 2025. \url{https://arxiv.org/abs/2508.15487}
\bibitem{qwen25} Qwen Team. \textit{Qwen2.5 Technical Report}. 2024. \url{https://arxiv.org/abs/2412.15115}
\bibitem{lavida} Li et al. \textit{LaViDa}. 2025. \url{https://arxiv.org/abs/2505.16839}
\bibitem{dreamvla} DreamLM. \textit{Dream-VL / Dream-VLA}. 2025--2026. \url{https://hkunlp.github.io/blog/2025/dream-vlx/}
\bibitem{lladavla} Wen et al. \textit{LLaDA-VLA}. 2025. \url{https://wenyuqing.github.io/llada-vla/}
\bibitem{discretediffvla} \textit{Discrete Diffusion VLA}. 2025. \url{https://arxiv.org/abs/2508.20072}
\bibitem{diffusiongemma_docs} Google. \textit{DiffusionGemma}. 2026. \url{https://ai.google.dev/gemma/docs/diffusiongemma}
\bibitem{crossway} Li et al. \textit{Crossway Diffusion}. ICRA 2024. \url{https://arxiv.org/abs/2307.01849}
\bibitem{oat} Liu et al. \textit{Ordered Action Tokens}. 2026. \url{https://arxiv.org/abs/2607.21670}
\bibitem{dp3} Ze et al. \textit{3D Diffusion Policy}. 2024. \url{https://arxiv.org/abs/2403.03954}
\bibitem{libero} Liu et al. \textit{LIBERO}. NeurIPS D\&B 2023. \url{https://arxiv.org/abs/2306.03310}
\end{thebibliography}

\end{document}
